\section{Distinguishing identified effects from unresolved comparisons in Java}
\label{exec19:javaid-scope}

A lending analyst may ask whether approving a loan changes repayment, or whether
a particular declined borrower would have repaid if approved. Those questions
differ from updating a supplied Bayesian network after observing borrower
attributes. They also differ from collecting regression coefficients over
possible graph orientations. A Java integration must preserve these distinctions
before it exposes a probability or an effect. Interface names alone cannot
determine the scientific contract.

The selected Tetrad source, version 7.6.11-SNAPSHOT at preserved revision
\texttt{b89afaf952a4}, contains an \texttt{Identifiability} interface in addition
to \texttt{Ida} and \texttt{PdagPagIda}. The identification implementation
constructs c-components and Q-expressions and evaluates them using an
observational joint-probability representation. It is a distinct candidate
from continuous-data effect enumeration. The source header names a 2002
algorithm report, but the report's proof has not been inspected for this
comparison. Code observations establish interface and failure-path requirements;
they do not establish theorem completeness or unrestricted latent-variable
identification. Nine selected original files support a bounded Tetrad comparison,
not an exhaustive census of Java software.~\citep{exec19_javaid_source}

\subsection{Observation, intervention and individual counterfactuals}
\label{exec19:javaid-witnesses}

Consider two authored lending mechanisms. Let $U$ be a binary borrower state,
$A$ approval, and $Y$ a binary repayment outcome. Let $U$ have probability
$1/2$ of being one, and let independent noise variables $N,E$ each have
probability $1/4$ of being one. Write $\oplus$ for exclusive-or and set
\begin{equation}
 A=U\oplus N,\qquad
 \mathcal M_1:\ Y=U,\qquad
 \mathcal M_2:\ Y=A\oplus E.
 \label{exec19:javaid-models}
\end{equation}
Both mechanisms are compatible with the allowed supergraph
$U\to A$, $U\to Y$, $A\to Y$ when an allowed arrow may have zero functional
effect. This construction does not impose faithfulness. Both yield the same
strictly positive observed joint law,
\begin{equation}
 \begin{array}{c|cc}
  &Y=0&Y=1\\\hline
 A=0&3/8&1/8\\
 A=1&1/8&3/8
 \end{array},
 \qquad P(Y=1\mid A=1)=3/4.
 \label{exec19:javaid-law}
\end{equation}
Nevertheless, the intervention risks are $(1/2,1/2)$ in $\mathcal M_1$ and
$(1/4,3/4)$ in $\mathcal M_2$, ordered by approval zero and one. Their average
causal effects are zero and $1/2$. The positive observed law does not select
one of these effects under the declared model class. These are two explicit
witnesses to nonuniqueness, not an enumeration of a sharp identified set. A
conditional probability cannot be relabeled as an intervention risk merely
because both are between zero and one.~\citep{exec19_javaid_reference}

For a purely illustrative monetary conversion, stipulate USD 200 per unit
increase in repayment probability and an approval-action cost of USD 20. The
two mechanisms give incremental net values
\begin{equation}
 200\{E[Y\mid\mathrm{do}(A=1)]-E[Y\mid\mathrm{do}(A=0)]\}-20
 \in\{-20,80\}\quad\text{USD per case}.
 \label{exec19:javaid-money}
\end{equation}
The sign changes without changing the observed joint law. These amounts are
stipulated arithmetic, not loan-pricing evidence, a measured bank return or
the valuation of a real approval policy. Additional causal restrictions or
intervention evidence would be needed to choose a mechanism for decision use.
\citep{exec19_javaid_reference}

Even marginal intervention risks may not settle an individual counterfactual.
In a second exact construction, let $A$ and $U$ be independent fair binary
variables. Compare potential outcomes $(Y(0),Y(1))=(U,U)$ with
$(Y(0),Y(1))=(U,1-U)$. Both generate the uniform observed distribution over
$(A,Y)$ and both have intervention risks $1/2$ for either approval value. For
an observed case with $A=0,Y=0$, however,
\begin{equation}
 P\{Y(1)=1\mid A=0,Y=0\}=
 \begin{cases}0,&(Y(0),Y(1))=(U,U),\\
 1,&(Y(0),Y(1))=(U,1-U).
 \end{cases}
 \label{exec19:javaid-coupling}
\end{equation}
The missing information concerns the joint coupling of potential outcomes,
which the two marginal intervention distributions do not supply. An agent
translating the business question must specify whether it asks for a population
intervention, conditional intervention or cross-world event. Substituting a
Bayesian conditional for the last answers another question.
\citep{exec19_javaid_reference}

\subsection{The selected updater has a narrower evidence contract}
\label{exec19:javaid-evidence}

The selected \texttt{BayesUpdater} documentation describes a joint marginal
conditional on evidence. The \texttt{Identifiability} implementation instead
comments its target as $P_t(s)$, with intervention values obtained from its
evidence object. \texttt{Evidence.getVariablesInEvidence} collects variables
with a single allowed category without checking manipulation flags.
\texttt{Identifiability.getJointMarginal} uses those variables as interventions
and explicitly assumes they are manipulated. Its separate graph-construction
method removes incoming edges only where manipulation flags are true.
Ordinary and manipulated evidence are therefore not interchangeable merely
because they share a container. A first proposed profile should admit validated
singleton interventions and reject mixed evidence until that mismatch is
resolved by a reviewed interface and native tests.~\citep{exec19_javaid_source}

There is also a subtype obligation. The public constructor accepts
\texttt{BayesIm}, but the no-evidence path and Q-expression evaluator cast it
to \texttt{MlBayesImObs} to access the observed joint-probability table. A
generic Bayesian-model object is insufficient evidence that the call will work.
The wrapper must preserve the observed-variable/category mapping and validate
the representation, finite probabilities, normalization and support. The
selected probability setter permits NaN; the Q-expression evaluator computes
a reciprocal denominator without a local zero guard. These observations
require numerical and support checks. They do not prove which end-to-end
native inputs produce a nonfinite result.~\citep{exec19_javaid_source}

Failures need distinct types. A recursive identification failure returns
\texttt{null}; the public joint-marginal path maps that branch to $-1$.
Internal consistency failures throw exceptions, whereas contradictory target
and intervention assignments return the valid probability zero and removing
all target variables returns one. The updated-model getter itself returns
\texttt{null}. The inspected GUI wrapper formats numeric marginals but does
not supply a typed identification certificate in that path. Clamping $-1$ to
zero would fabricate a probability; treating a null model as a zero effect
would confuse object availability with causal meaning. Signed treatment
contrasts are a separate type and can legitimately be negative, so a universal
negative-means-failure rule would also be wrong.~\citep{exec19_javaid_source}

\subsection{An effect list is not a confidence set}
\label{exec19:javaid-candidates}

The IDA-related interfaces need separate graph and sampling contracts.
The \texttt{PdagPagIda} documentation mentions PAGs, but its constructor checks
for a legal DAG or PDAG and continuous data. A nearby comment says that the
rPAG-IDA route is disabled because it still needs work. A class name, dormant
branch or broader comment cannot admit a PAG-only input. Fully oriented graphs
may satisfy more than one graph predicate; the observation is the actual guard,
not a categorical claim about every object labeled a PAG.
\citep{exec19_javaid_source}

Candidate failures complicate interpretation. In a selected optimal-adjustment
branch, an exception can append zero and continue. An outer handler instead
logs an error and omits a candidate. A returned zero or shortened list cannot
automatically certify a structural zero or exhaustive graph-compatible effect
enumeration. A proposed adapter should retain every attempted candidate with
its graph, adjustment set, regression status and exclusion reason, and refuse
scientific publication when failures collapse into values. This source-level
observation is neither a native reproduction nor a claim that every call fails.
\citep{exec19_javaid_source}

The distance functions show another distinction. For candidates $\{-2,2\}$
and target zero, finite-set distance is two, whereas distance to the interval
hull $[-2,2]$ is zero. The selected code uses the extreme range for multiple
candidates. For an empty list it returns the target value itself, so target
minus one yields minus one rather than a valid nonnegative distance. These
source-branch examples require an explicit empty status and an accurately named
metric. Neither an effect list nor its hull is a calibrated confidence region
without an argument for graph uncertainty, sampling error and observation law.
\citep{exec19_javaid_source,exec19_javaid_reference}

\subsection{A proposal route with explicit unresolved obligations}

The proposed Java schema separates graph and probability inputs, singleton
interventions, ordinary evidence, query variables and cross-world conditioning.
Its output distinguishes a completed identified query, a reviewed
nonidentification certificate, unknown scope, unsupported input, and execution
or numerical failure. The source's $-1$ branch must not be upgraded to a
universal nonidentification certificate before the algorithm's assumptions and
proof are reviewed. The original-language engine remains a candidate; no port,
native qualification or distribution permission follows from this comparison.

The remaining public research must inspect the named algorithm original,
graph-legality and probability-storage helpers, and other Java candidates under
a recorded search boundary. It must establish whether conditional ID, general
counterfactual or partial-identification queries have actual supported
interfaces, rather than infer absence from selected files. A later admitted
native prototype can test subtypes, manipulation flags, failure sentinels, zero
support, nonfinite values, skipped candidates and graph guards against the
exact references. R and Python symbolic engines are distinct comparators with
their own graph and failure contracts. These tasks can use public sources and
synthetic laws without requiring authentic bank data to clarify the proposal.
